\begin{abstract}
Large Language Model (LLM)-based software engineering tools are commonly
evaluated by their task performance, while their runtime, token usage, and
failure behavior receive less attention. We extend CASCADE~\cite{Kiecker_2026},
a tool for detecting inconsistencies between source code and documentation, with a configurable,
run-scoped metrics recording architecture. The instrumentation captures
structured events for pipeline timings, LLM attempts and token consumption,
repair activity, Java and Docker execution, and failures. It writes a detailed
JSON Lines event stream and an aggregated JSON summary without replacing
CASCADE's analysis result or suppressing its errors. We use this extension to
compare three locally hosted LLMs and one externally hosted LLM on the balanced
142-sample core benchmark. The results show that no model dominates across all
criteria: gpt-oss~\cite{ollama_gpt_oss_20b} achieves the strongest detection
result, qwen2.5-coder~\cite{ollama_qwen2.5-coder-14b-instruct} the lowest mean
runtime, and mistral-small~\cite{ollama_mistral-small3.2_24b} the lowest mean
token usage. However, the
high failure rates limit the reliability of the detection comparison and prevent
attributing the observed differences to model compatibility alone. The recorder
therefore provides a basis for examining efficiency and reliability trade-offs
alongside CASCADE's detection metrics.

\keywords{code--documentation consistency \and LLM evaluation \and metrics
instrumentation \and CASCADE \and software engineering}
\end{abstract}
